% ---------------------------------------------------------------------
% III. METHOD (about 0.6 page)
% Owners: one paragraph per experiment owner: ___
%
% FACTS lists below were checked against the notebooks (*.ipynb), the
% results CSVs and the journal. Source in [brackets] after each fact.
% "nb" = notebook. Where a fact could NOT be confirmed it says CHECK.
% ---------------------------------------------------------------------
\section{Method}\label{sec:method}

\subsection{Filters}
% PLAN: which filters did we test?
% - Detoxify "original" and Detoxify "unbiased".
% - Note: they are different model types (BERT vs RoBERTa), not just
%   debiased vs not debiased.
% - Our own simple, transparent filter built for Exp 4 (a logistic
%   regression on words and word pairs), and how accurate it is compared
%   with Detoxify (numbers in Exp 4's exp4_auc.csv).
%
% FACTS:
% - Detoxify models: "original" and "unbiased" [all nbs, Cell 1]
% - Output used: the "toxicity" score, P(toxic) [all nbs, proba()]
% - Converted to log-odds for SHAP: log(p / (1 - p)), p clipped to
%   1e-6 .. 1 - 1e-6 [all nbs, logit()]
% - Mask tokens: "[MASK]" for original, "<mask>" for unbiased
%   [Exp 3 nb, Cell 9] -> fits BERT vs RoBERTa
% - BERT vs RoBERTa: [journal, Dixon 2018 note] -- CHECK against the
%   Detoxify README (cite key: detoxify); not stated in the notebooks
% - "unbiased" was trained for the Jigsaw Unintended Bias competition,
%   i.e. on Civil Comments [Exp 4 nb, Cell 4 caveat; journal, Borkan 2019
%   note] -> cite borkan2019nuanced
% - Transparent filter [Exp 4 nb, Cells 2-3]:
%   - dataset: Civil Comments (google/civil_comments, Hugging Face)
%   - toxic = toxicity >= 0.5
%   - training: 100,000 comments sampled from the train split (seed 42)
%   - features: binary word counts, words that appear in >= 5 comments
%   - main filter: words + word pairs (unigrams + bigrams)
%   - sanity-check filter: single words only
%   - model: logistic regression (C = 1.0)
% - Accuracy, AUC on 5,000 test comments (seed 42) [exp4_auc.csv]:
%   - transparent, words + pairs: 0.892
%   - transparent, words only: 0.897
%   - Detoxify original: 0.928
%   - Detoxify unbiased: 0.971 (favoured: trained on this dataset)

\todo{Detoxify original and unbiased, plus the transparent filter (Exp 4).}

\subsection{Sentence Sets}
% PLAN: what sentences did we feed the filters?
% - Exp 1: 48 sentences that use the same insult in three ways
%   (promoting, reporting, arguing against).
% - Exp 2: 190 harmless sentences where only one word changes
%   (an identity word like "gay", or a neutral word like "tall").
%   Based on Dixon 2018.
% - Exp 4 also uses 200 real comments from Civil Comments.
% - The other experiments reuse these same sentences.
%
% FACTS:
% - Exp 1 set: 48 sentences = 16 topics x 3 versions [Exp 1 nb, Cell 2;
%   exp1_results.csv has 48 rows]
%   - promote / report / counter, same insult words in all 3
%   - example (topic 1): "Those fans are disgusting idiots and should be
%     banned." / "The club condemned a post calling its fans disgusting
%     idiots who should be banned." / "Calling the fans disgusting idiots
%     is unfair and wrong."
%   - report sentences are longer than the others [Exp 1 nb, Cell 2
%     comment] -> limitation
% - Exp 2 set: 190 sentences = 10 harmless templates x 19 words
%   [Exp 2 nb, Cell 2; exp2_results.csv has 190 rows]
%   - 15 identity words: gay, lesbian, transgender, black, white, asian,
%     muslim, jewish, christian, immigrant, deaf, blind, disabled, poor,
%     rural
%   - 4 control words: tall, young, local, quiet
%   - 150 identity sentences + 40 control sentences
%   - example template: "My best friend is {}."
%   - design follows Dixon 2018 [journal, Dixon 2018 note] -> cite
%     dixon2018measuring
% - Exp 4 set: 438 sentences [exp4_sentences.csv]
%   - the 48 Exp 1 + 190 Exp 2 sentences
%   - plus 200 real comments: 100 toxic, 100 non-toxic, 5-14 words,
%     from the full Civil Comments test split (seed 42) [Exp 4 nb, Cell 6]
%   - cap of 16 tokens; all 438 were under it, none removed
% - Which experiment uses which set [nbs, Cell 2 of each]:
%   - Exp 3: the 190 Exp 2 sentences
%   - Exp 5 and Exp 6: all 238 Exp 1 + Exp 2 sentences

\todo{Exp 1, Exp 2, and the Exp 4 real comments.}

\subsection{Explainers}
% PLAN: how did we set up SHAP and LIME?
% - SHAP: which version we used and its main setting.
% - LIME: how many samples and that we fixed the random seed.
% - Both explain each word's effect on the "toxic" score.
% - When an explainer "removes" a word, something has to replace it.
%   Say what we used, and that Exp 3 tests other options.
% - Careful: the settings are not the same in every experiment
%   (Exps 1-2 used fewer samples than Exps 3-6). Check the notebooks.
%
% FACTS:
% - All run on Google Colab, T4 GPU; libraries detoxify, shap, lime
%   [all nbs, Cell 0-1]; random seed 42 [all nbs]
% - Units: whole words, keeping attached punctuation, so SHAP and LIME
%   score the same words [Exps 1, 3-6 nbs, word_tokenizer()]
%   - Exp 2 split on spaces instead [Exp 2 nb, Cell 1]; Exp 3 re-ran
%     Exp 2's sentences with the word-level version [Exp 3 nb]
% - SHAP [all nbs]:
%   - shap.Explainer with a text masker -> Partition explainer (Exp 4
%     names "partition" explicitly)
%   - explains log-odds of toxic (so values add up)
%   - max_evals = 500 (Exp 2: 300)
%   - efficiency check after every Detoxify run: base value + sum of
%     SHAP values must equal the log-odds within 0.01, or the notebook
%     stops [all nbs, check_efficiency()]
%   - deterministic: Exp 2 compares two SHAP runs [Exp 2 nb, Cell 8]
%     -- CHECK: the printed result was not saved in the notebook
% - LIME [all nbs, lime_word_weights()]:
%   - LimeTextExplainer, labels = (1,) so weights are for "toxic"
%   - explains the probability, not log-odds -> only compared by rank
%     with SHAP
%   - one feature per distinct word, so no word is cut off
%   - samples and seeds per experiment:
%     - Exp 1: 2000 samples, seed 42
%     - Exp 2: 2000 samples, seed 42 (agreement on 40 sentences); 500
%       and 5000 samples x seeds 0-9 (stability tests)
%     - Exp 3, Exp 5: 5000 samples, seed 42
%     - Exp 4: 5000 samples, seeds 0-9
%     - Exp 6: 500 and 5000 samples, seeds 0-9
% - Removal baseline (what replaces a "removed" word):
%   - SHAP default: "..." [Exp 3 nb, Cell 9 note says "..." is shap's
%     default and what Exps 1-2 used] -- CHECK: the printed value was not
%     saved in the notebook
%   - Exp 3 also tests: deleting the word (a single space, because an
%     empty string crashes shap) and the model's mask token ("[MASK] " /
%     "<mask> ") [Exp 3 nb, Cell 9]
%   - Exp 4: deletion, for both SHAP and the exact values [Exp 4 nb]
%   - Exp 5's deletion test deletes words [Exp 5 nb, drop()]
% - Which model each experiment explains:
%   - both Detoxify models: Exps 1-3 SHAP, Exp 3 LIME, Exp 5
%   - original only: Exp 1-2 LIME, Exp 6
%   - transparent filter only: Exp 4

\todo{SHAP and LIME configuration, log-odds, removal baseline.}

\subsection{References}
% PLAN: what "right answers" did we compare the explainers against?
% - Counterfactual effect (Exp 3): swap the identity word for a neutral
%   word and see how much the score changes. Only one word differs, so
%   that change is the word's real effect.
% - Exact Shapley values (Exp 4): on our transparent filter we can work
%   out the exact correct answer, then see if SHAP and LIME match it.
% - Why we need these: two explainers agreeing doesn't prove they are
%   right (Krishna 2022).
%
% FACTS:
% - Counterfactual effect [Exp 3 nb, Cell 5]:
%   - for each of the 150 identity sentences
%   - = its log-odds minus the mean log-odds of the 4 control sentences
%     from the same template
%   - SHAP's estimate to compare = the identity word's SHAP value minus
%     the mean SHAP value of the control words in the same template
%   - "spill-over" = counterfactual effect minus SHAP's estimate (the
%     part SHAP put on other words)
%   - for LIME: the same, but in probability [Exp 3 nb, Cell 6]
%   - counterfactual swap method -> cite garg2019counterfactual
%     (in refs.bib; no journal review yet)
% - Exact Shapley values [Exp 4 nb, Cell 7]:
%   - on the transparent filter only
%   - players = the same word tokens SHAP uses
%   - value of a subset = the filter's log-odds with all other words
%     deleted
%   - every subset is tried (2^n, n <= 16) and the Shapley formula is
%     applied exactly
%   - efficiency holds exactly (checked to 1e-6)
%   - compared against: SHAP Partition, SHAP Permutation (both deletion,
%     max_evals = 500) and LIME (10 seeds, 5000 samples) [Exp 4 nb, Cell 8]
% - Sanity checks [Exp 4 nb, Cell 10]:
%   - on the single-word filter, SHAP Partition should match exactly
%   - words with no weight (null players) should get 0
% - Why references are needed -> cite krishna2022disagreement [journal,
%   Krishna 2022 note]

\todo{Counterfactual effect (Exp 3) and exact Shapley values (Exp 4).}

\subsection{Metrics}
% PLAN: what did we measure?
% - Faithfulness (Exps 3-4): how closely the explainers match the
%   right answers.
% - Deletion test (Exp 5): delete the words the explainer blames and see
%   if the score drops more than when deleting random words.
% - Stability (Exp 6): run LIME several times and see if the answer
%   changes.
% - Signal strength (Exp 6): how strong the toxic signal is in each
%   sentence, to see if weak signals give worse explanations.
%
% FACTS:
% - Faithfulness to the counterfactual effect [Exp 3 nb, faith()]:
%   - Pearson r, Spearman rho, sign agreement, mean absolute error (MAE,
%     log-odds)
%   - at sentence level (n = 150) and term level (n = 15)
%   - LIME: Spearman and sign agreement only (probability units)
% - Faithfulness to exact Shapley [Exp 4 nb, Cells 8-9]:
%   - per sentence: Spearman rho, top-word match, MAE (SHAP only)
%   - split into low- vs high-interaction sentences (interaction = how
%     much the word pairs change the exact values; split at the median)
% - Removal baseline effect [Exp 3 nb, Cell 9]:
%   - faithfulness metrics repeated for each of the 3 baselines
%   - Spearman rho between baselines' rankings of the 15 identity words
% - Deletion test [Exp 5 nb, Cells 4-5], following ERASER -> cite
%   deyoung2020eraser:
%   - top k = 1, 2, 3 words from SHAP, from LIME, and k random words
%     (10 random draws, averaged)
%   - comprehensiveness = log-odds drop when those words are deleted
%     (higher = better)
%   - sufficiency = log-odds drop when only those words are kept
%     (lower = better)
%   - AOPC = average over k = 1, 2, 3
%   - 238 sentences x 2 models x 3 k values [exp5_deletion_by_sentence.csv]
% - Stability [Exp 6 nb, Cells 4-5]:
%   - 238 sentences x 2 sample sizes x 10 seeds = 4,760 LIME runs
%     [exp6_lime_runs.csv]
%   - flip rate = share of runs whose top word differs from the most
%     common top word
%   - rank stability = mean Spearman rho between every pair of runs
%   - also: mean LIME R squared (fit), agreement with SHAP's top word
% - Signal strength [Exp 6 nb, Cell 3]:
%   - |log-odds - SHAP base value| (how far the score is from the score
%     with every word removed)
% - Warning signs [Exp 6 nb, Cell 6]:
%   - on the 150 identity sentences, 5000-sample LIME runs
%   - error = |counterfactual effect - SHAP estimate| (from Exp 3)
%   - tested: low R squared, SHAP-LIME top-word disagreement, weak signal
%   - measured by Spearman rho with the error and by AUC for flagging the
%     worst 25% of errors

\todo{Faithfulness, deletion, stability, signal strength.}
